% Präambel aller Bände; erwartet \Sprache (german oder ngerman) und \Titel
\documentclass[
  paper=5.5in:8.5in,
  BCOR=16mm,
  DIV=calc,
  fontsize=11pt,
  twoside,
  open=right,
  headings=small,
  footnotes=multiple,
  cleardoublepage=empty,
  numbers=noenddot]{scrbook}
% Zeilenabstand 13 pt, Fußnoten 11,14 pt: 7 Fußnotenzeilen auf 6 Textzeilen
\linespread{0.9559}\selectfont
% fehlende Glyphen als Fehler
\tracinglostchars=3
\addtokomafont{footnote}{\linespread{1.013}\selectfont}
\usepackage[polutonikogreek,main=\Sprache,shorthands=off]{babel}
\babelprovide[hyphenrules=\Sprache-x-latest]{\Sprache}
% Schreibung vor 1901 („Werth“, „Theil“): nicht zwischen t und h trennen
\makeatletter
\directlua{lang.patterns(lang.new(\number\l@german), "t4h")}
\makeatother
\usepackage{fontspec}
\directlua{
  luaotfload.add_fallback("kapital",
    {"file:GentiumBook-Regular.ttf:mode=node;"})}
% Variable Schrift: Grundschrift 400, Fußnoten und kleiner etwas kräftiger
\setmainfont{CrimsonPro-VF.ttf}[
  ItalicFont=CrimsonPro-Italic-VF.ttf,
  BoldFont=CrimsonPro-VF.ttf,
  BoldItalicFont=CrimsonPro-Italic-VF.ttf,
  Ligatures={Common,TeX},
  Numbers={OldStyle,Proportional},
  RawFeature={fallback=kapital},
  UprightFeatures={SizeFeatures={
    {Size=-10, RawFeature={axis={wght=460}}},
    {Size=10-, RawFeature={axis={wght=400}}}}},
  ItalicFeatures={SizeFeatures={
    {Size=-10, RawFeature={axis={wght=460}}},
    {Size=10-, RawFeature={axis={wght=400}}}}},
  BoldFeatures={SizeFeatures={{Size=-, RawFeature={axis={wght=650}}}}},
  BoldItalicFeatures={SizeFeatures={{Size=-, RawFeature={axis={wght=650}}}}}]
\newfontfamily\greekfont{GentiumBook-Regular.ttf}[
  ItalicFont=GentiumBook-Italic.ttf,
  Scale=MatchLowercase]
% DIV=calc erst mit Grundschrift und Zeilenabstand
\normalfont\recalctypearea
\catcode`℔=\active \protected\def℔{\lbsign}
\catcode`‥=\active \protected\def‥{.\kern.15em.}
\newcommand{\lbsign}{\leavevmode\hbox{\setbox0\hbox{lb}\copy0\kern-\wd0
  \raise1.08ex\hbox to \wd0{%
    \hss\vrule height .065ex depth 0pt width 1.1\wd0\hss}}}
\usepackage{realscripts}
\usepackage{amsmath}
\usepackage{graphicx}
\usepackage{array}
\usepackage{booktabs}
\usepackage{xcolor}
\usepackage{enumitem}
\usepackage{scrlayer-scrpage}
\usepackage[
  protrusion=true,
  expansion=true]{microtype}
% Erster Band für den Druck in zwei Teilen: Seitenzahlen des anderen Teils
\ifdefined\KapitalOther
  \usepackage{xr-hyper}
  \externaldocument[x-]{\KapitalOther}
\fi
\usepackage{refcount}
\usepackage{hyperref}
\usepackage{bookmark}
\urlstyle{same}

\setkomafont{pagehead}{\normalfont\itshape\small}

\setlength{\parindent}{1em}
\clubpenalty=3000
\widowpenalty=10000
\displaywidowpenalty=10000
\raggedbottom
\setlength{\parskip}{0pt}
\setlist[itemize]{leftmargin=1em, labelwidth=1em, labelsep=0pt, align=left,
  topsep=0pt, partopsep=0pt, itemsep=0pt, parsep=0pt}
\let\KapitalFootnoterule\footnoterule
\renewcommand*{\footnoterule}{\vfill\KapitalFootnoterule}
% Marke grundlinienbündig, rechtsbündig vor dem Text; der Einzug richtet sich
% nach der längsten Nummer des Bandes
\newlength{\fnindent}
\deffootnote{\fnindent}{1em}{\thefootnotemark\enspace}
\setkomafont{footnotelabel}{\addfontfeatures{Numbers={OldStyle,Monospaced}}}
\providecommand{\FnWidest}{00}
\makeatletter
\AtBeginDocument{%
  \settowidth{\fnindent}{\footnotesize\ftnm@font\FnWidest\enspace}}
\makeatother
% Crimson Pro hat hochgestellte Ziffern, aber keine Buchstaben: das „a“ in „19a“
% verkleinert und angehoben wie die Ziffern
\makeatletter
\deffootnotemark{\fnsup{\thefootnotemark}}
\newcommand{\fnsup}[1]{\edef\kfn{#1}\let\kfnd\@empty\let\kfnl\@empty
  \expandafter\@tfor\expandafter\x\expandafter:\expandafter=\kfn\do{%
    \ifcat a\x\edef\kfnl{\kfnl\x}\else\edef\kfnd{\kfnd\x}\fi}%
  \textsuperscript{\kfnd}%
  \ifx\kfnl\@empty\else
    \raisebox{.19em}{\fontsize{\fpeval{.67*\f@size}}{0}\selectfont\kfnl}%
  \fi}
\makeatother
\setlength{\skip\footins}{13pt plus 4pt minus 2pt}
% Fußnoten ohne Zwischenraum: Strut von 0,7 Fußnotenzeilen
\setlength{\footnotesep}{7.8pt}
\renewcommand{\arraystretch}{1.1}
\setlength{\emergencystretch}{3em}
\tolerance=1000
\makeatletter\@beginparpenalty=\@M\makeatother

\hypersetup{
  hidelinks,
  bookmarksnumbered=false,
  pdfpagemode=UseOutlines,
  pdfauthor={Karl Marx},
  pdftitle={\Titel},
  pdflang={de}}

% Fußnote mit der Marke des Drucks
\newcommand{\fn}[2]{{\renewcommand{\thefootnote}{#1}\footnote{#2}}}
% dasselbe für Tabellen: Marke in der Zelle, Text hinter der Tabelle
\newcommand{\fnmark}[1]{{\renewcommand{\thefootnote}{#1}\footnotemark}}
\newcommand{\fntext}[2]{{\renewcommand{\thefootnote}{#1}\footnotetext{#2}}}
% Seitenwechsel der Erstausgabe: Strich im Text, Seitenzahl am äußeren Rand;
% \leavevmode, damit die Marke am Absatzanfang nicht auf der Vorseite landet
\newcommand{\pb}[2]{\leavevmode\phantomsection\label{#1}%
  \textcolor{black!45}{\raisebox{-.15ex}{\rule{.3pt}{1.9ex}}}\pbnum{#2}}
\newcommand{\pbm}[2]{\leavevmode\phantomsection\label{#1}\pbnum{#2}}
% Ziffer per \vadjust in Höhe der Zeile, ohne \marginpar: dessen Ausgaberoutine
% bringt geteilte Fußnoten durcheinander; die Seite kommt aus \Ifthispageodd
\newcommand{\pbfont}{\footnotesize\color{black!60}}
\newcommand{\pbnum}[1]{\vadjust{\kern-\dp\strutbox
  \vbox to 0pt{\boxmaxdepth=0pt\vss\Ifthispageodd
    {\hbox to 0pt{\hskip\dimexpr\textwidth+\marginparsep\relax\pbfont #1\hss}}%
    {\hbox to 0pt{\hss\pbfont #1\hskip\marginparsep}}}%
  \kern\dp\strutbox}}
\newcommand{\pbx}[1]{\phantomsection\label{#1}}

% Überschriften; die Kolumnentitel setzt der Konverter
\newcommand{\hdI}[1]{\par\vspace*{3\baselineskip}%
  {\centering\normalfont\Large\addfontfeatures{LetterSpace=-0.5}#1\par}%
  \nobreak\vspace{2\baselineskip}\nobreak}
\newcommand{\hdII}[1]{\par\vspace{2\baselineskip plus 1\baselineskip}%
  {\centering\normalfont\large #1\par}%
  \nobreak\vspace{1.2\baselineskip}\nobreak}
\newcommand{\hdIII}[1]{\par\vspace{1.5\baselineskip plus .5\baselineskip}%
  {\centering\normalfont #1\par}%
  \nobreak\vspace{.8\baselineskip}\nobreak}
\newcommand{\hdIV}[1]{\par\vspace{1\baselineskip plus .5\baselineskip}%
  {\centering\normalfont\itshape #1\par}%
  \nobreak\vspace{.5\baselineskip}\nobreak}
\newcommand{\hdV}[1]{\par\vspace{.8\baselineskip plus .3\baselineskip}%
  {\centering\normalfont\itshape #1\par}%
  \nobreak\vspace{.3\baselineskip}\nobreak}
\let\hdVI\hdV

% Blockformen; sie behalten einen umgebenden Einzug (Fußnoten), den \centering
% und \raggedright auf null setzen
\makeatletter
\newcommand{\kcentering}{%
  \skip@=\leftskip\centering\leftskip=\dimexpr\skip@\relax plus 1fil\relax}
\newcommand{\kraggedright}{\skip@=\leftskip\raggedright\leftskip=\skip@}
\makeatother
% zentrierte Zeilen sind meist Tabellen- oder Formelüberschriften: nicht von der
% Folgezeile trennen (@nobreak gilt auch für die Ausgaberoutine einer
% Randziffer in der Zeile)
\makeatletter
\newcommand{\cenline}[1]{\par{\kcentering\global\@nobreaktrue #1\par}%
  \nobreak\global\@nobreakfalse}
\makeatother
% Formelzeile, die nicht umbrechen darf, notfalls auf Satzbreite verkleinert
\newcommand{\fitline}[1]{\par\noindent\setbox0\hbox{\displayfractrue
  \addfontfeatures{Numbers=Lining}#1}%
  \ifdim\wd0>\linewidth\resizebox{\linewidth}{!}{\box0}\else\box0\fi\par}
\newcommand{\dispf}[1]{\par\smallskip{\kcentering\displayfractrue
  \addfontfeatures{Numbers=Lining}#1\par}\smallskip}
\newcommand{\etline}[1]{\par{\kraggedright\advance\leftskip2em\noindent #1\par}}
\newcommand{\rline}[1]{\par{\raggedleft #1\par}}
\newenvironment{hanglist}[1][1.5em]
  {\par\smallskip\begingroup\kraggedright\advance\leftskip#1\relax}
  {\par\endgroup\smallskip}
\newcommand{\hitem}[1][1em]{\hangindent=#1\hangafter=1\noindent}
\newenvironment{bquote}
  {\par\smallskip\begingroup\advance\leftskip2em\parindent=0pt}
  {\par\endgroup\smallskip}
% Zu breite Tabellen ragen höchstens 3 mm in den Bundsteg, sonst nach außen bis
% 6 mm vor den Rand, notfalls verkleinert
\newsavebox\wtbox
\newdimen\wtover
\newdimen\wtin
\newdimen\wtroom
\newcommand{\widetab}[1]{\leavevmode\sbox\wtbox{#1}%
  \wtover=\dimexpr\wd\wtbox-\linewidth\relax
  \ifdim\wtover>0pt
    \wtin=.5\wtover \ifdim\wtin>3mm \wtin=3mm\fi
    \wtroom=\dimexpr\paperwidth-1in-\oddsidemargin-\textwidth-6mm\relax
    \ifdim\dimexpr\wtover-\wtin\relax>\wtroom
      \sbox\wtbox{%
        \resizebox{\dimexpr\linewidth+\wtin+\wtroom\relax}{!}{\usebox\wtbox}}%
    \fi
    \Ifthispageodd
      {\makebox[\linewidth][l]{\hspace*{-\wtin}\usebox\wtbox}}%
      {\makebox[\linewidth][r]{\usebox\wtbox\hspace*{-\wtin}}}%
  \else\makebox[\linewidth][c]{\usebox\wtbox}\fi}
\newcommand{\tabfont}{%
  \addfontfeatures{Numbers={Lining,Monospaced}}\setlength{\tabcolsep}{.4em}}
\newcommand{\linerule}[1]{\kern.12em\rule[.3em]{#1}{.4pt}\kern.12em}
% Brüche im Fließtext eine Stufe kleiner und mit 2 pt Luft zu den Nachbarzeilen,
% in abgesetzten Formeln in voller Größe; Versalziffern, damit Zähler und Nenner
% gleichmäßig am Bruchstrich stehen
\newif\ifdisplayfrac
\newcommand{\Frac}[2]{\text{\addfontfeatures{Numbers=Lining}%
  \ifdisplayfrac$\dfrac{#1}{#2}$\else
    \raisebox{0pt}[\dimexpr\height+2pt][\dimexpr\depth+2pt]{\small$\dfrac{#1}{#2}$}%
  \fi}}
\newcommand{\ovarc}[1]{\overline{#1}}
\newcommand{\lab}[2]{\overbrace{#2}^{\text{\scriptsize #1}}}
\newbox\brzhalf
\newcommand{\nt}[2]{\overset{\text{\scriptsize #1}}{#2}}
\newcommand{\vfrac}[2]{{\addfontfeatures{Fractions=On}#1⁄#2}}
\newcommand{\grc}[1]{\foreignlanguage{polutonikogreek}{\greekfont\upshape #1}}
\providecommand{\TeilVermerk}{}
\newcommand{\maintitlefont}{\Huge\addfontfeatures{LetterSpace=-1.2}}
\newcommand{\authorfont}{\LARGE\addfontfeatures{LetterSpace=-0.8}}

% Ursprüngliches Inhaltsverzeichnis mit den Seiten dieser Ausgabe
\newcommand{\tocpage}[1]{%
  \ifcsname r@#1\endcsname\pageref{#1}\else\pageref{x-#1}\fi}
\newcommand{\tocline}[3]{\par{%
  \leftskip=#1\relax\rightskip=2.5em plus 1fil\parfillskip=-\rightskip
  \parindent=0pt\hangindent=1em\noindent #2%
  \nobreak\leaders\hbox to .6em{\hss.\hss}\hfill\nobreak
  \hbox to 2.5em{\hss\tocpage{#3}}\par}}
\newcommand{\tocbook}[1]{\par\medskip{\centering\itshape #1\par}\smallskip}
